\section{Meta-Reinforcement Learning}


\textbf{Setup:} Please navigate to \texttt{src/submission/meta\_rl} and follow the instructions in the README.

\textbf{Code Overview:} The main entry point for the code is via \texttt{src/submission/meta\_rl/scripts/dream.py} and \texttt{src/submission/meta\_rl/scripts/rl2.py}, which are the training scripts for DREAM and RL$^2$ respectively. Both of these can be invoked as follows:

\begin{lstlisting}[language=bash]
$ python -m meta_rl.scripts.{script} {experiment_name} -b environment=\"map\"
\end{lstlisting}

In this invocation, \texttt{\{script\}} can either be \texttt{dream} or \texttt{rl2} and \texttt{\{experiment\_name\}} can be any string with no white spaces. Results from this invocation are saved under \texttt{experiments/\{experiment\_name\}}. For example, to launch the DREAM training script and save the results to \texttt{experiments/dream}, you would run:

\begin{lstlisting}[language=bash]
$ python -m meta_rl.scripts.dream dream -b environment =\"map\"
\end{lstlisting}

To overwrite previous results, you would run:

You can pass the \texttt{--force\_overwrite} or \texttt{-f} flag to run another experiment with the same experiment name, which will overwrite any previously saved files at the corresponding experiment directory. An example command for running DREAM is below:

\begin{lstlisting}[language=bash]
$ python -m meta_rl.scripts.dream dream -b environment=\"map\" -f
\end{lstlisting}

If you do not pass this flag, the scripts will not allow you to run two experiments with the same experiment name.

There are two important subdirectories in each experiment directory:

\begin{itemize}
    \item \textbf{Tensorboard:} Each experiment includes a \texttt{experiments/experiment\_name/tensorboard} subdirectory, which will log important statistics about the training run, including the meta-testing returns under the \texttt{rewards/test} tag and the meta-training returns under the \texttt{rewards/training} tag. To view these, point Tensorboard at the appropriate directories. All curves are plotted with two versions, one where the x-axis is number of meta-training trials under \texttt{tensorboard/episode} and one where the x-axis is the number of environment steps \texttt{tensorboard/step}.
    
    \item \textbf{Visualizations:} Each experiment also includes a \texttt{experiments/experiment\_name/visualize} subdirectory. This directory includes \texttt{.gif} videos of the agent during meta-testing and is structured as follows. The top level of subdirectories identify how many meta-training trials have elapsed before the video.

    Note that in the visualizations saved under \texttt{experiments/experiment\_name/visualize}:

    \begin{itemize}
        \item The agent is rendered as a red square.
        
        \item The grid cells that the agent has visited in the episode are rendered as small origin squares.
        
        \item There are four pairs of buses, rendered as blue, pink, cyan, and yellow squares.
        
        \item The map is rendered as a black square.
        
        \item The goal state is rendered as a green square, which obscures one of the buses.
    \end{itemize}

    
    In experiments run with \texttt{dream.py}, the exploration episode is saved under \texttt{\{video\_num\}-exploration.gif} and the exploration episode is saved under \texttt{\{video\_num\}-exploitation.gif}. For example, the video under \texttt{experiments/dream/visualize/10000/0-exploration.gif} is the first exploration meta-testing episode after 10000 meta-training trials.
    
    In experiments run with \texttt{rl2.py}, \texttt{\{video\_num\}.gif} contains both the exploration and exploitation episodes, with the exploration episode first. For example, the video under \texttt{experiments/rl2/visualize/10000/0.gif} contains the first exploration and exploitation episode after 10000 meta-training trials.
\end{itemize}

You will implement two short methods inside the \texttt{embed/encoder\_decoder.py} file.

\begin{enumerate}[(a)]
    \input{03-meta-rl/01-grid-world-navigation}
    \input{03-meta-rl/02-e2e-meta-rl}
    \item \points{2c} \textbf{DREAM}

In part (a), we observed some shortcomings of end-to-end and existing decoupled meta-RL approaches. In this problem, we'll implement some components of DREAM, which attempts to address these shortcomings, given the assumption that each meta-training task is assigned a unique \emph{problem ID} $\mu$. During meta-testing, DREAM does not assume access to the problem ID.

From a high level, DREAM works by separately learning exploration and exploitation. To learn exploitation, DREAM learns an exploitation policy $\pi_\phi^{\text{exp}}(a \mid s, z)$ that tries to maximize returns during exploitation episodes, conditioned on a task encoding $z$. DREAM learns a \emph{encoder} $F_\psi(z \mid \mu)$ to produce the task encoding $z$ from the problem ID $\mu$. Critically, this encoder is trained in such a way that $z$ contains only the information necessary for the exploitation policy $\pi_\phi^{\text{exp}}$ to solve the task and achieve high returns.

By training the encoder in this way, DREAM can then learn to explore by trying to recover the information contained in $z$. To achieve this, DREAM learns an exploration policy $\pi_\phi^{\text{exr}}$, which produces an exploration trajectory $\tau^{\text{exp}} = (s_0, a_0, r_0, \ldots)$ when rolled out during the exploration episode. To recover the information contained in $z$, DREAM fits a maximize the mutual information between the encoding $z$ and the exploration trajectory $\tau^{\text{exp}}$:
\[
\max I(F_\psi(z \mid \mu), \tau^{\text{exp}}).
\]

This objective can be optimized by maximizing the following variational lower bound:
\[
\mathcal{J}(\omega, \phi) = \mathbb{E}_{\mu \sim P_\mu, \tau^{\text{exp}} \sim \pi_\phi^{\text{exr}}}[\log q_\omega(z \mid \tau^{\text{exp}})]
\]
where $q_\omega(z \mid \tau^{\text{exp}})$ is a learned \emph{decoder}. Note that this decoder enables us to convert an exploration trajectory $\tau^{\text{exp}}$ into a task encoding $z$ that the exploitation policy uses. This is critical for meta-test time, where the problem ID is unavailable, and $z$ cannot be computed via the encoder $F_\psi(z \mid \mu)$.

The objective $\mathcal{J}(\omega, \phi)$ is optimized with respect to both the decoder $q_\omega$ and the exploration policy $\pi_\phi^{\text{exr}}$.


For simplicity, we parametrize the decoder $q_\omega(z \mid \tau^{\text{exp}})$ as a Gaussian $\mathcal{N}(g_\omega(\tau^{\text{exp}}), \sigma^2 I)$ centered at a learned $g_\omega(\tau^{\text{exp}})$ with unit variance. Then, $\log q_\omega(z \mid \tau^{\text{exp}})$ equals negative mean-squared error $-\|g_\omega(\tau^{\text{exp}}) - \text{stop-gradient}(z)\|_2^2$ plus some constants independent of $\omega$. Overall, \emph{maximizing} $\mathcal{J}(\omega, \phi)$ with respect to the decoder parameters $\omega$ is equal to \emph{minimizing} the below mean-squared error \textbf{with respect to} $\omega$:
\[
\mathbb{E}_{z \sim F_\psi(\mu)} \left[\|g_\omega(\tau^{\text{exp}}) - \text{stop-gradient}(z)\|_2^2\right].
\]
    
\textbf{Code:} Fill in the \texttt{.compute\_losses} method of the \texttt{EncoderDecoder} in \texttt{encoder\_decoder.py} to implement the above equation for optimize $\mathcal{J}(\omega, \phi)$ with respect to the decoder parameters $\omega$.


    \item \points{2d} To optimize $\mathcal{J}(\omega, \phi)$ with respect to the exploration policy parameters $\phi$, we expand out $\mathcal{J}(\omega, \phi)$ as a telescoping series:
\[
\mathcal{J}(\omega, \phi) = \mathbb{E}_{\mu \sim P_\mu(\mu)}[\mathcal{J}[g_\omega(z \mid s_0)]] + \mathbb{E}_{\mu \sim P_\mu(\mu), \tau^{\text{exp}} \sim \pi_\phi^{\text{exr}}} \left[\sum_{t=0}^{T-1} \log q_\omega(z \mid \tau_{t+1}^{\text{exp}}) - \log q_\omega(z \mid \tau_t^{\text{exp}})\right],
\]
where $\tau_t^{\text{exp}}$ denotes the exploration trajectory up to timestep $t$: $(s_0, a_0, r_0, \ldots, s_t)$. Only the second term depends on the exploration policy, and since it occurs per timestep, it can be maximized by treating it as the following intrinsic reward function $r_t^{\text{int}}$, which we can maximize with standard reinforcement learning:
\begin{equation}
r_t^{\text{int}}(a_t, r_t, s_{t+1}, \tau_{t+1}^{\text{exr}}, \mu) = \mathbb{E}_{z \sim F_\psi(\mu)} \left[\log q_\omega(z \mid \tau_{t+1}^{\text{exr}}) - \log q_\omega(z \mid \tau_t^{\text{exp}})\right].
\end{equation}

Note that $\tau_{t+1}^{\text{exr}}$ is equal to $\tau_t^{\text{exp}}$ with the additional observations of $(a_t, r_t, s_{t+1})$. The exploration policy is optimized to maximize this intrinsic reward $r_t^{\text{int}}$ instead of the extrinsic rewards $r_t$, which will maximize the objective $\mathcal{J}(\omega, \phi)$. Intuitively, $r_t^{\text{int}}$ is the ``information gain'' representing how much additional information about $z$ -- which encodes all the information to solve the task -- the tuple $s_{t+1}, r_t, a_t, s_t$ contains over what was already observed in $\tau_t^{\text{exp}}$.

\textbf{Code:} Implement the reward function $r_t^{\text{int}}(a_t, r_t, s_{t+1}, \tau_{t+1}^{\text{exr}}, \mu)$ by filling in the \texttt{label\_rewards} function of \texttt{EncoderDecoder} in \texttt{encoder\_decoder.py}. Note that you'll need to make the same substitution for $\log q_\omega(z \mid \tau^{\text{exp}})$ in Equation (1) that we used in part a).

    \input{03-meta-rl/05-dream-analyze-returns}
    \input{03-meta-rl/06-dream-analyze-behavior}
\end{enumerate}